% !TEX program = xelatex
% =====================================================================
%  全国大学生数学建模竞赛（CUMCM）论文 LaTeX 模板 —— 自包含单文件
%  文件：main.tex（正文）+ refs.bib（参考文献）——只需要这两个文件，
%        不需要额外的 preamble/宏文件，直接把两个文件拷进你的工作目录即可。
% ---------------------------------------------------------------------
%  编译（Windows，MiKTeX / TeX Live 均可）：
%      xelatex main.tex        % 第 1 遍
%      bibtex  main            % 只在参考文献变动时需要
%      xelatex main.tex        % 第 2 遍：解析交叉引用
%      xelatex main.tex        % 第 3 遍：让 \ref 与页码稳定
%  一行版：latexmk -xelatex main.tex
%  若你暂时不想用 .bib，可把 \bibliography{refs} 换成手写 thebibliography
%  环境（见正文“参考文献”处的注释）。
%
%  中文字体：ctexart + fontset=windows 直接调用 Windows 自带的宋体/黑体/
%  楷体，Windows 上开箱即用。Linux/macOS 请把 fontset=windows 换成
%  fontset=fandol（随 TeX 发行版自带）。详见同目录 README.md。
% =====================================================================
%
%  【匿名合规 —— 国赛硬性要求】
%      《论文格式规范（2026 年修订稿）》要求全文匿名：摘要页、正文、附录
%      任何地方都不能出现参赛者姓名、学校、赛区信息。
%      本模板已移除作者/学校/队号字段（没有 \author、没有队号变量），
%      请勿自行加回；也不要使用带学校名的页眉页脚或水印。
%
%  【规则年份标注 —— 逐条以当年官方通知为准】
%      以下条款均来自 references/contests.md 记录的官方文件，规则逐年修订，
%      提交前必须重新核对当年官方通知（国赛官网 https://www.mcm.edu.cn/ ）：
%      (a) 电子版第 1 页必须是摘要专用页，承诺书与编号专用页不放入电子版
%          ——《论文格式规范（2026 年修订稿）》
%      (b) 摘要（含标题与关键词）原则上不超过一页；页码从摘要页开始，
%          页脚中部、阿拉伯数字连续编号——同上（本模板页脚已按此设置）
%      (c) 正文不要目录、不超过 30 页；附录页数不限——同上
%      (d) 附录必须含【支撑材料文件列表】+【建模用到的全部完整、可运行
%          源程序】；缺少必要源程序、程序不能运行或运行结果与论文不符，
%          都可能被取消评奖资格——同上
%      (e) 四边页边距 ≥2.5 cm，左侧装订——同上（见下方 geometry 设置）
%      (f) AI 工具使用声明必须排在【参考文献之前】——《AI 工具使用规定
%          （2026 年试行）》
%      (g) 全国规范不统一规定字体字号，各赛区可另作要求 —— 定稿前务必
%          向本赛区确认字号字体，不要照搬网上某个学校的模板
% =====================================================================

\documentclass[12pt,a4paper,fontset=windows]{ctexart}

% ------------------------- 版面与字体 -------------------------
\usepackage{geometry}
\geometry{a4paper,left=2.5cm,right=2.5cm,top=2.5cm,bottom=2.5cm} % 四边 ≥2.5cm
\usepackage{setspace}
\usepackage{indentfirst}          % 每节首段也缩进（中文排版习惯）

% ------------------------- 数学 -------------------------
\usepackage{amsmath}
\usepackage{amssymb}
\usepackage{bm}

% ------------------------- 表格与图 -------------------------
\usepackage{booktabs}             % 三线表 \toprule \midrule \bottomrule
\usepackage{multirow}
\usepackage{array}
\usepackage{tabularx}
\usepackage{graphicx}
\usepackage{float}                % 提供 [H] 浮动体位置
\usepackage{caption}
\usepackage{subcaption}
\usepackage{tikz}                 % 技术路线图（流程图）自己画，不用外部图片
\usetikzlibrary{arrows.meta,positioning,shapes.geometric,calc}
\usepackage{pgfplots}             % 演示用图；你自己的图建议用 matplotlib 导出 PDF
\pgfplotsset{compat=1.18}

% ------------------------- 算法伪代码 -------------------------
\usepackage{algorithm}
\usepackage{algpseudocode}        % algorithmicx 的 pseudocode 语言
\floatname{algorithm}{算法}
\renewcommand{\algorithmicrequire}{\textbf{输入：}}
\renewcommand{\algorithmicensure}{\textbf{输出：}}
\renewcommand{\algorithmicend}{\textbf{结束}}
\renewcommand{\algorithmicif}{\textbf{如果}}
\renewcommand{\algorithmicthen}{\textbf{则}}
\renewcommand{\algorithmicelse}{\textbf{否则}}
\renewcommand{\algorithmicfor}{\textbf{对}}
\renewcommand{\algorithmicforall}{\textbf{对所有}}
\renewcommand{\algorithmicdo}{\textbf{执行}}
\renewcommand{\algorithmicwhile}{\textbf{当}}
\renewcommand{\algorithmicloop}{\textbf{循环}}
\renewcommand{\algorithmicrepeat}{\textbf{重复}}
\renewcommand{\algorithmicuntil}{\textbf{直到}}
\renewcommand{\algorithmicreturn}{\textbf{返回}}
\renewcommand{\algorithmicprocedure}{\textbf{过程}}
\renewcommand{\algorithmicfunction}{\textbf{函数}}

% ------------------------- 代码清单（Python 高亮） -------------------------
\usepackage{listings}
\usepackage{xcolor}
\definecolor{codekw}{RGB}{0,0,180}
\definecolor{codecmt}{RGB}{0,128,0}
\definecolor{codestr}{RGB}{160,32,32}
\definecolor{codebg}{RGB}{248,248,248}
\lstset{
  language=Python,
  basicstyle=\ttfamily\small,        % 中文注释跟着 xeCJK 走，用 \small 保证不溢出
  keywordstyle=\color{codekw}\bfseries,
  commentstyle=\color{codecmt},
  stringstyle=\color{codestr},
  numbers=left,
  numberstyle=\tiny\color{gray},
  numbersep=6pt,
  frame=single,
  rulecolor=\color{gray!50},
  backgroundcolor=\color{codebg},
  breaklines=true,                   % 长行自动折行，避免溢出版心
  breakatwhitespace=false,
  postbreak=\mbox{\textcolor{gray}{$\hookrightarrow$}\space},
  extendedchars=true,                % 允许 UTF-8 / 中文出现在代码与注释里
  inputencoding=utf8,
  columns=flexible,
  keepspaces=true,
  showstringspaces=false,
  tabsize=4,
  captionpos=t,
  xleftmargin=1.5em,
  aboveskip=1em,belowskip=1em
}

% ------------------------- 页眉页脚（国赛：只用页脚页码） -------------------------
\usepackage{fancyhdr}
\pagestyle{fancy}
\fancyhf{}
\fancyfoot[C]{\thepage}            % 页码：页脚中部，从摘要页开始连续编号
\renewcommand{\headrulewidth}{0pt}
\renewcommand{\footrulewidth}{0pt}

% ------------------------- 标题样式与交叉引用 -------------------------
\captionsetup{font=small,labelsep=quad}      % 图题在图下、表题在表上（LaTeX 默认）
\ctexset{
  section/format   = {\heiti\zihao{4}\raggedright},
  subsection/format= {\heiti\zihao{-4}\raggedright},
  subsubsection/format = {\heiti\zihao{-4}\raggedright}
}
\usepackage[hidelinks]{hyperref}             % 便于电子版跳转；不需要可删

% =====================================================================
\begin{document}

% =====================================================================
% 第 1 页：摘要专用页（电子版第一页，不含承诺书与编号专用页）
% 提示：摘要的权重很高——评委常常先读摘要再决定是否细读正文，
%       并且摘要不合格会被直接降档。摘要【必须含量化结果】（误差、
%       最优值、提升比例、灵敏度结论），不能只写“建立了合理的模型”。
%       按正文定稿后再回来写摘要（写作顺序见 assets/abstract-template.md）。
%       本页不要出现学校/姓名/队号。
% =====================================================================
\begin{center}
  {\zihao{3}\heiti 【论文题目：具体、含研究对象与方法，避免“关于某问题的研究”】}
\end{center}

\vspace{0.6em}
\begin{center}
  {\zihao{4}\heiti 摘\quad 要}
\end{center}
\vspace{0.3em}

% 摘要正文：建议 600–1000 字，每一问都要“方法 + 具体数值结果 + 检验结论”。
针对问题一，本文建立【模型类型，如多目标混合整数规划】模型。首先对【数据/条件】进行【预处理/分析】，
得到【关键量，含单位】；其次以【目标函数】为目标、以【约束条件】为约束，构建【模型名】；
然后采用【求解算法/工具及版本】求解，得到【具体数值结果，含单位】。最后以【检验方式】检验，
【误差/偏差指标】为【数值】，表明模型【结论】。

针对问题二，在问题一的基础上引入【新增因素/机制】，建立【模型名】。由【方法】求解得到【结果】，
与【基线方案：现状/贪心/单目标模型】相比，【关键指标】改善【百分比】。对【关键参数】做灵敏度
分析表明，当【参数】在【范围】内变动时，【目标量】的波动不超过【数值】，模型具有较好的稳健性。

针对问题三，建立【模型名】并设计【策略/方案】，求解得到【核心结论】，据此给出【可执行建议】。

本文的创新点在于：【1–3 条真创新，如新增约束、改进算法、组合建模、数据融合】。

\vspace{0.5em}
\noindent{\heiti 关键词：}\ 【关键词1】；【关键词2】；【关键词3】；【关键词4】

\newpage
% =====================================================================
% 正文（不要目录；正文不超过 30 页——2026 年修订稿）
% =====================================================================

\section{问题重述}
用自己的语言复述背景、已知条件与待求目标，\textbf{不要直接复制题面}。

\subsection{问题背景}
【背景、数据来源、实际意义；数据来源须可追溯，网上资料要进参考文献】
% 正文必须出现 \cite 标注，否则参考文献列表是空的（bibtex 会报
% “I found no \citation commands”）。示例：
% 本文的建模思路参考了文献~\cite{jiang2018model,si2021algorithm}。
本文的建模思路与求解框架参考了文献~\cite{jiang2018model,si2021algorithm}。

\subsection{待解决的问题}
\begin{itemize}
  \item 问题一：【用一句话说清已知什么、要求什么、输出形式是什么（数值/方案/策略）】
  \item 问题二：【……】
  \item 问题三：【……】
\end{itemize}

\section{问题分析}
每问都要说清三件事：核心矛盾是什么、关键变量是哪些、为什么选这个模型而不选替代方案。

\subsection{问题一的分析}
【核心矛盾 + 关键变量 + 可行方法对比（例如“最短路模型 vs 整数规划模型”，说明取舍理由）】
问题含离散与连续两类变量，凸优化框架下的混合整数规划可以精确刻画该结构~\cite{boyd2004convex}；
若需要同时优化多个目标，可用多目标进化算法~\cite{deb2002nsga2} 求 Pareto 前沿。

\subsection{问题二的分析}
【在问题一结论上的增量分析；新引入的机制及其必要性】

\subsection{问题三的分析}
【……】

\subsection{技术路线图}
% 硬要求：问题分析这一节必须有一张技术路线图——评委靠它 10 秒判断你是否想清楚。
% 下面这张图用 TikZ 直接画在 .tex 里，因此模板不依赖任何外部图片文件。
\begin{figure}[htbp]
  \centering
  \begin{tikzpicture}[
      every node/.style={font=\small},
      box/.style={draw, rounded corners=2pt, align=center, minimum height=8mm,
                  inner xsep=5pt, fill=blue!5},
      ar/.style={-{Stealth[length=2mm]}, thick}
    ]
    \node[box] (data) {数据预处理\\与指标筛选};
    \node[box, right=12mm of data] (q1) {问题一：\\【模型名】};
    \node[box, right=12mm of q1]  (q2) {问题二：\\【模型名】};
    \node[box, right=12mm of q2]  (q3) {问题三：\\【模型名】};
    \node[box, below=10mm of data] (val) {误差分析与\\模型检验};
    \node[box, below=10mm of q2]   (sens){灵敏度分析\\与稳健性讨论};
    \node[box, below=10mm of q3]   (dec) {结论与\\可执行建议};
    \draw[ar] (data) -- (q1);
    \draw[ar] (q1) -- (q2);
    \draw[ar] (q2) -- (q3);
    \draw[ar] (q1) -- (val);
    \draw[ar] (q2) -- (sens);
    \draw[ar] (q3) -- (dec);
  \end{tikzpicture}
  \caption{本文总体技术路线图（请替换为与你的方案一致的流程）}
  \label{fig:roadmap}
\end{figure}

如图~\ref{fig:roadmap} 所示，本文按“数据预处理 $\rightarrow$ 分问建模求解 $\rightarrow$ 检验与灵敏度”的主线展开。

\section{模型假设}
假设建议 3--8 条；每条都要写“内容 + 理由 + 对模型的影响”，只有内容的假设会被认为不严谨。
\begin{enumerate}
  \item \textbf{假设 1}：【内容】。理由：【依据，尽量给出数据或文献支撑】。对模型的影响：【影响】。
  \item \textbf{假设 2}：【内容】。理由：【……】。对模型的影响：【……】。
  \item \textbf{假设 3}：【内容】。理由：【……】。对模型的影响：【……】。
\end{enumerate}

\section{符号说明}
% 经验数据（社区对 2021–2025 年 64 篇国赛获奖论文的统计，非官方）显示：
% “符号说明”出现率约 94%，比摘要还高——中文建模论文的事实标准，默认就写。
\begin{table}[htbp]
  \centering
  \caption{主要符号说明}
  \label{tab:notation}
  \begin{tabular}{clc}
    \toprule
    符号 & 含义 & 单位 \\
    \midrule
    $x_{ij}$   & 决策变量：【含义】 & 【单位】 \\
    $c_{ij}$   & 【含义】           & 【单位】 \\
    $Z$        & 目标函数值：【含义】 & 【单位】 \\
    $\lambda$  & 【含义，如权重系数】 & 无量纲 \\
    \bottomrule
  \end{tabular}
\end{table}

表~\ref{tab:notation} 中未列出的符号在首次出现处给出定义；全文术语与符号必须保持一致。
多指标权重可用层次分析法标定~\cite{saaty1980ahp}，权重来源与一致性检验结果应在此处说明。

\section{模型的建立与求解}
主体章节（建议占正文 10--15 页）。每一问都按“为什么选这个模型 $\rightarrow$ 模型建立（公式）
$\rightarrow$ 求解算法 $\rightarrow$ 结果（图/表）”四段式写。

\subsection{问题一：【模型名称】}

\subsubsection{为什么选这个模型}
【动机 + 与替代方案的对比。例：问题含离散选址与连续分配两类变量，线性规划无法表达
“选址后才有流量”的逻辑，因此选混合整数规划；相比启发式，小规模下可求全局最优。】

\subsubsection{模型建立}
决策变量、目标函数与约束如下：
\begin{align}
  \min\ Z & = \sum_{i \in V}\sum_{j \in J} c_{ij} x_{ij}
             + \lambda \sum_{j \in J} f_j y_j
             \label{eq:objective} \\[2pt]
  \text{s.t.}\quad
  & \sum_{j \in J} x_{ij} = d_i, \qquad \forall i \in V, \label{eq:demand} \\[2pt]
  & \sum_{i \in V} x_{ij} \le M y_j, \qquad \forall j \in J, \label{eq:capacity} \\[2pt]
  & x_{ij} \ge 0,\quad y_j \in \{0,1\}. \label{eq:domain}
\end{align}
其中式~\eqref{eq:objective} 为目标函数，式~\eqref{eq:demand} 保证需求被满足，
式~\eqref{eq:capacity} 刻画“只有建站才能服务”的逻辑（$M$ 为足够大的常数），
式~\eqref{eq:domain} 为取值范围。模型适用条件与局限：【一句话说清前提，例：需求确定为已知】。

\subsubsection{求解算法设计}
\begin{algorithm}[H]
  \caption{【算法名，如：拉格朗日松弛 + 局部搜索】求解问题一}
  \label{alg:q1}
  \begin{algorithmic}[1]
    \Require 需求集合 $\{d_i\}$、候选站集合 $J$、距离矩阵 $(c_{ij})$、最大迭代次数 $K$
    \Ensure 最优选址方案 $y^{*}$ 及目标函数值 $Z^{*}$
    \State 初始化：$y^{(0)} \gets$ 【初始解构造方式】，$Z^{(0)} \gets +\infty$
    \For{$k = 1$ \textbf{到} $K$}
      \State 固定 $y^{(k-1)}$，解式~\eqref{eq:demand}--\eqref{eq:capacity} 得到流量分配 $x^{(k)}$
      \State 按【邻域动作，如 1-opt 交换】生成候选解集 $\mathcal{N}(y^{(k-1)})$
      \If{$\exists\, y' \in \mathcal{N}(y^{(k-1)})$ 使 $Z(y') < Z^{(k-1)}$}
        \State $y^{(k)} \gets \arg\min_{y' \in \mathcal{N}} Z(y')$，更新 $Z^{(k)}$
      \Else
        \State $y^{(k)} \gets y^{(k-1)}$；\textbf{跳出}循环
      \EndIf
    \EndFor
    \State \Return $y^{(K)},\, Z^{(K)}$
  \end{algorithmic}
\end{algorithm}

算法~\ref{alg:q1} 的时间复杂度为【$O(\cdot)$】；参数设置为【步长/惩罚系数/迭代上限】，
收敛判据为【判据】。局部搜索部分可参考无导数直接搜索方法~\cite{nelder1965simplex}。
软件与版本：【Python 3.11 + numpy 1.26 / LINGO 20 / MATLAB R2023b】。

\subsubsection{求解结果}
\begin{figure}[htbp]
  \centering
  \begin{tikzpicture}
    \begin{axis}[
        width=0.72\textwidth, height=0.40\textwidth,
        xlabel={迭代次数 $k$},
        ylabel={目标函数值 $Z$（万元）},
        grid=major,
        legend pos=north east,
        legend style={font=\small, draw=none, fill=none},
        every axis plot/.append style={thick}
      ]
      % 【占位数据】下面两组数值仅用于演示排版，请替换为你自己算出的结果序列。
      \addplot[mark=*, color=blue] coordinates
        {(0,128.4)(5,112.7)(10,104.2)(15,101.3)(20,100.2)(30,100.0)};
      \addlegendentry{本文算法}
      \addplot[mark=square*, dashed, color=red] coordinates
        {(0,128.4)(5,120.1)(10,115.6)(15,113.2)(20,112.4)(30,112.0)};
      \addlegendentry{贪心基线}
    \end{axis}
  \end{tikzpicture}
  \caption{算法收敛曲线与基线对比（横轴：迭代次数；纵轴：目标函数值，单位万元）}
  \label{fig:convergence}
\end{figure}

\begin{table}[htbp]
  \centering
  \caption{问题一求解结果（【占位数据】，请替换为程序输出）}
  \label{tab:q1result}
  \begin{tabular}{lcc}
    \toprule
    方案 & 目标函数值 $Z$ / 万元 & 服务覆盖率 / \% \\
    \midrule
    现状方案（基线） & 128.4 & 86.5 \\
    贪心算法       & 112.0 & 91.2 \\
    本文模型       & 100.0 & 96.8 \\
    \bottomrule
  \end{tabular}
\end{table}

结果如图~\ref{fig:convergence} 与表~\ref{tab:q1result} 所示：本文算法在第【$k$】次迭代后收敛，
目标函数值较基线下降【$x\%$】。所有数值均由附录程序输出，避免手抄数字造成论文与代码不符。

\subsection{问题二：【模型名称】}
【同问题一的四段式；重点写清在问题一基础上的增量，以及新增变量的物理意义】

\subsection{问题三：【模型名称】}
【同前】

\section{结果分析与检验}
% 检验与灵敏度不一定要独立成章（社区统计：约 23%/19% 的获奖论文独立成章），
% 但内容必须存在——只给结果不做检验是明确的失分点。

\subsection{误差分析}
【误差指标定义（RMSE/MAE/MAPE）、计算方式、数值；解释“为什么误差是这个量级”
（数据精度、模型简化、随机性），不要只报一个数】

\subsection{与基准方案的对比}
【至少与一个基线对比：现状方案 / 贪心 / 单目标模型 / 文献结果。给出表格与百分比变化，
说明改进来自哪一步（模型结构？求解质量？数据融合？）】

\subsection{灵敏度分析}
【选 2--4 个对结论影响最大的参数；扰动范围给依据（$\pm 10\%/\pm 20\%$，或按数据精度）。
结论句模板：“当【参数】在【范围】内变动时，【目标量】变化不超过【$Z\%$】，模型稳健”
或“……高度敏感，实际中需精确测量该参数”。】
全局灵敏度指标（如 Sobol 指数）的定义与使用可参考~\cite{saltelli2008global}。

\subsection{模型检验}
【用独立数据/留出法验证；或做极端情形测试（边界值、退化情形）；或与解析解/小规模
穷举解比对。说明检验通过与否，不通过时如何修正】

\section{模型的评价与推广}
\subsection{模型优点}
【3 条以上，每条对应一个可验证的事实（如“可求全局最优”“复杂度为 $O(n^2)$”“对缺失数据稳健”）】
\begin{enumerate}
  \item 【优点 1】
  \item 【优点 2】
  \item 【优点 3】
\end{enumerate}

\subsection{模型缺点与改进方向}
【2 条以上，写成“局限 + 改进方向”，诚实但不自我否定】
\begin{enumerate}
  \item 【局限 1】$\rightarrow$ 改进方向：【……】
  \item 【局限 2】$\rightarrow$ 改进方向：【……】
\end{enumerate}

\subsection{模型推广}
【适用于什么条件、可迁移到哪些场景（写清楚前提，不要空喊“可广泛应用”）】

% =====================================================================
% AI 工具使用声明
% 位置硬要求：《AI 工具使用规定（2026 年试行）》规定须在【参考文献之前】
% 设置本声明。未使用与已使用二选一，照抄对应那句即可，不要两边都留。
% 注意：故意隐瞒 AI 使用、虚假声明，或把未经人工核验的 AI 输出直接作为
%       核心建模与分析成果提交的，取消评奖资格。
%       使用 AI 的，支撑材料中须含《AI工具使用详情.pdf》，内容包括：工具
%       名称/版本、使用目的与环节、主要提示方式与过程、对输出的采纳与
%       人工修改核验情况。
% =====================================================================
\section*{AI 工具使用声明}
\addcontentsline{toc}{section}{AI 工具使用声明}   % 本模板无目录，此行可留可删

% —— 情况一：未使用任何 AI 工具 ——
本参赛队在竞赛过程中未使用任何AI工具。

% —— 情况二：使用了 AI 工具（把上面那句注释掉，改用下面这句并填写）——
% 本参赛队在竞赛过程中使用了AI工具，主要用于【简要用途，如语言润色、代码调试等】，
% 详细使用情况见支撑材料《AI工具使用详情.pdf》。

% =====================================================================
% 参考文献
% 国赛规范：引用他人或公开资料（含网上资料）须按科技论文规范列出，并在正文标注。
% 用 .bib：把 refs.bib 与本文件放同一目录，编译时跑一次 bibtex main。
% 不用 .bib：把下面 \bibliography{refs} 两行换成手写列表亦可，格式：
%   \begin{thebibliography}{99}
%     \bibitem{ref1} 作者. 题名. 出处, 年份, 卷(期): 页码.
%   \end{thebibliography}
% =====================================================================
\bibliographystyle{unsrt}          % 按正文引用顺序编号 [1][2][3]
\bibliography{refs}

\newpage
% =====================================================================
% 附录（页数不限，须与正文装订；电子版中同样不得出现身份信息）
% 硬要求：附录必须包含【支撑材料文件列表】+【建模用到的全部完整、可运行
% 源程序】（含 Excel/SPSS 等交互命令）。缺少必要源程序、程序不能运行或
% 运行结果与论文不符，都可能被取消评奖资格。确实未用程序时写明
% “本论文没有用到程序”。
% =====================================================================
\appendix
\section{支撑材料文件列表}
% 支撑材料：全部可运行源程序 + 自采数据 + 大篇幅中间结果图表，压成一个
% RAR/ZIP（≤20 MB），文件列表必须放进论文附录；无则写“本论文没有支撑材料”。
\begin{table}[htbp]
  \centering
  \caption{支撑材料文件列表（【占位内容】，请按实际情况增删）}
  \label{tab:files}
  \begin{tabular}{llp{7.2cm}}
    \toprule
    文件 & 类型 & 说明 \\
    \midrule
    \texttt{code/q1\_solve.py} & Python 源码 & 问题一建模与求解，输出 \texttt{results/q1.json} \\
    \texttt{code/q2\_solve.py} & Python 源码 & 问题二求解 \\
    \texttt{code/plot\_fig.py} & Python 源码 & 生成论文中的全部图 \\
    \texttt{data/raw.csv}      & 数据       & 原始数据（来源：【给出出处】） \\
    \texttt{results/}          & 结果        & 程序输出的 JSON/CSV，论文数值均来自此处 \\
    \bottomrule
  \end{tabular}
\end{table}

\section{主要源程序}
% 推荐用 \lstinputlisting 直接引入源文件，保证论文代码与实际运行的代码完全一致：
%   \lstinputlisting[language=Python, caption={问题一求解程序 q1_solve.py}]{code/q1_solve.py}
% 下面的 lstlisting 只是一段可在本文件内直接编译的示例（去掉它换成你的代码）。

\begin{lstlisting}[caption={问题一求解程序（示例，请替换为你的完整代码）}, label={lst:q1}]
# -*- coding: utf-8 -*-
"""问题一求解程序：最小二乘拟合并输出结果（示例）

运行环境：Python 3.9+，仅依赖 numpy（与论文结果一一对应）。
输出文件：results.json —— 论文表格/图中的数值全部来自该文件，
          避免手抄数字导致“论文数字与附录代码不符”。
"""

import json

import numpy as np


def fit_poly(x, y, degree):
    """最小二乘多项式拟合。

    参数:
        x: 自变量序列
        y: 因变量序列
        degree: 多项式阶数

    返回:
        多项式系数（按次数升序）
    """
    design = np.vander(np.asarray(x, dtype=float), degree + 1)
    coef, *_ = np.linalg.lstsq(design, np.asarray(y, dtype=float), rcond=None)
    return coef


def main():
    rng = np.random.default_rng(20260910)      # 固定种子，保证结果可复现
    x = np.linspace(0.0, 10.0, 60)
    y = 2.0 + 0.8 * x - 0.05 * x ** 2 + rng.normal(0.0, 0.1, x.size)
    coef = fit_poly(x, y, 2)
    y_hat = np.vander(x, 3) @ coef
    rmse = float(np.sqrt(np.mean((y - y_hat) ** 2)))
    with open("results.json", "w", encoding="utf-8") as fp:
        json.dump({"coef": coef.tolist(), "rmse": rmse}, fp,
                  ensure_ascii=False, indent=2)
    print("RMSE = %.4f" % rmse)


if __name__ == "__main__":
    main()
\end{lstlisting}

程序~\ref{lst:q1} 为示例结构，完整程序见支撑材料 \texttt{code/} 目录；
若本论文确实没有用到程序，请在此处写明“本论文没有用到程序”。

\end{document}
